\section{Đánh giá}
\label{sec:evaluation}

\subsection{Bộ dữ liệu}
\label{subsec:datasets}

Bộ dữ liệu nền chứa các sự kiện lành tính tổng hợp và ba loại rủi ro có nhãn:
SQL injection, đoán mật khẩu và truy cập trái phép. Bộ dữ liệu rủi ro theo ngữ
cảnh chứa thêm các kịch bản liên quan đến địa chỉ IP mới, hoạt động ngoài giờ,
truy cập tài nguyên hiếm gặp, hành động bị từ chối lặp lại và chuỗi tấn công
nhiều sự kiện. Bảng~\ref{tab:dataset-description} tách rõ hai tệp CSV và vai
trò của chúng. Trường \texttt{expected\_label} là nhãn phục vụ đánh giá trong
nguyên mẫu; nó không phải dữ liệu mà một triển khai vận hành sẽ có sẵn.

\begin{table}[t]
  \centering
  \small
  \caption{Các tập dữ liệu tổng hợp và vai trò đánh giá. Phân bố nhãn được đếm từ hai tệp CSV nguồn; tổng event/alert được đối chiếu với metadata artifact CP3.}
  \label{tab:dataset-description}
  \begin{tabular}{p{2.0cm}p{2.25cm}p{2.3cm}p{1.7cm}}
    \toprule
    Tập dữ liệu & Quy mô và nhãn & Vai trò & Kết quả tái sinh \\
    \midrule
    Baseline (\texttt{demo}) & 60 event: 30 benign, 10 SQL injection, 10 password guessing, 10 unauthorized access & Đánh giá detection nền theo nhãn & 20 alert \\
    Context (\texttt{risk-demo}) & 14 event: 4 benign, 2 context anomaly, 4 password guessing, 4 unauthorized access & Đánh giá theo kịch bản tín hiệu context và grouping & 15 alert, 10 finding, 3 incident \\
    \bottomrule
  \end{tabular}
\end{table}


\subsection{Thước đo}
\label{subsec:metrics}

Với bộ dữ liệu phát hiện theo luật nền, đánh giá báo cáo dương tính đúng (TP),
dương tính giả (FP), âm tính giả (FN), âm tính đúng (TN), độ chính xác
(precision), độ bao phủ (recall) và F1 theo từng loại rủi ro. Với lớp ngữ cảnh,
đánh giá được thực hiện theo kịch bản: số event hợp lệ, alert, finding theo
từng tín hiệu và incident được lấy trực tiếp từ artifact, thay vì gán một metric
phân loại độc lập cho từng finding. Cả hai lần chạy dùng cùng seed RBAC
\texttt{data/rbac\_seed.json} và \texttt{config/default.toml}: cửa sổ 300 giây,
giờ làm việc 08:00--18:00, và bonus ngữ cảnh 10 (IP mới), 8 (ngoài giờ/tài
nguyên hiếm), 12 (từ chối lặp lại) và 15 (chuỗi incident). Các lệnh tái sinh và
artifact đầu ra được lưu trong \texttt{PROGRESS.md} và thư mục \texttt{artifacts/}.

\subsection{Kết quả}
\label{subsec:results}

Bảng~\ref{tab:baseline-results} tóm tắt kết quả nền hiện tại.

\begin{table}[t]
  \centering
  \caption{Kiểm chứng các luật trên 60 event nhật ký tổng hợp. Các giá trị được tái sinh từ \texttt{artifacts/cp3-baseline/metrics.json}; chúng không đại diện cho hiệu năng vận hành.}
  \label{tab:baseline-results}
  \begin{tabular}{lrrrrrrr}
    \toprule
    Loại rủi ro & TP & FP & FN & TN & Precision & Recall & F1 \\
    \midrule
    SQL injection & 9 & 0 & 1 & 50 & 1.0000 & 0.9000 & 0.9474 \\
    Đoán mật khẩu & 10 & 0 & 0 & 50 & 1.0000 & 1.0000 & 1.0000 \\
    Truy cập trái phép & 10 & 0 & 0 & 50 & 1.0000 & 1.0000 & 1.0000 \\
    \midrule
    Trung bình macro & \multicolumn{4}{c}{--} & 1.0000 & 0.9667 & 0.9825 \\
    \bottomrule
  \end{tabular}
\end{table}


Lần chạy tái lập xử lý 60 sự kiện hợp lệ, không có dòng lỗi và tạo 20 cảnh
báo. Ba nhãn có precision tuyệt đối trong tập tổng hợp này; riêng SQL injection
có một âm tính giả, nên recall là 0.9000 và F1 là 0.9474. Do đó, các số liệu chỉ
cho thấy hành vi của nguyên mẫu trên kịch bản có nhãn, không phải hiệu năng trên
nhật ký vận hành thực tế.

Trên kịch bản rủi ro theo ngữ cảnh, lần chạy xử lý 14 sự kiện hợp lệ, không có
dòng lỗi, tạo 15 cảnh báo và 10 phát hiện ngữ cảnh. Bảng~\ref{tab:context-findings}
cho thấy sáu phát hiện ngoài giờ, hai phát hiện tài nguyên hiếm, một IP mới và
một cụm từ chối lặp lại. Số lượng finding không phải số sự kiện duy nhất: một
sự kiện có thể mang nhiều tín hiệu; vì vậy kết quả được trình bày như tín hiệu
triage giải thích được thay vì một độ đo phân loại độc lập.

\begin{table}[t]
  \centering
  \caption{Artifact của phần mở rộng context-risk trên 14 event tổng hợp. Số lượng lấy từ \texttt{artifacts/cp3-context/context\_findings.json}; đây không phải metric phân loại.}
  \label{tab:context-findings}
  \begin{tabular}{lrrr}
    \toprule
    Tín hiệu ngữ cảnh & Số finding & Confidence & Context bonus \\
    \midrule
    Ngoài giờ (CTX-AFTER-HOURS) & 6 & 0.75 & 8 \\
    IP mới (CTX-NEW-IP) & 1 & 0.80 & 10 \\
    Tài nguyên hiếm (CTX-RARE-RESOURCE) & 2 & 0.70 & 8 \\
    Từ chối lặp lại (CTX-REPEATED-DENIAL) & 1 & 0.85 & 12 \\
    \midrule
    Tổng & 10 & -- & -- \\
    \bottomrule
  \end{tabular}
\end{table}


Cuối cùng, cơ chế grouping tạo ba incident (Bảng~\ref{tab:incident-summary}).
Mỗi incident được tạo bằng cách nhóm cảnh báo theo cùng cặp người dùng--IP;
trong từng nhóm, một đoạn mới bắt đầu khi chênh lệch với cảnh báo trước lớn hơn
300 giây. Mức điểm incident là điểm cảnh báo lớn nhất trong đoạn cộng bonus chuỗi
khi đoạn có nhiều cảnh báo, sau đó bị chặn ở 100. Vì artifact chỉ lưu hai mốc
đầu-cuối, bảng không diễn giải nó như một timeline pháp chứng đầy đủ.

\begin{table*}[t]
  \centering
  \small
  \caption{Tóm tắt incident tái sinh từ \texttt{artifacts/cp3-context/incidents.json}. Thời gian là hai mốc \texttt{start\_time} và \texttt{end\_time} do chương trình ghi nhận, không phải timeline chi tiết.}
  \label{tab:incident-summary}
  \resizebox{\textwidth}{!}{%
  \begin{tabular}{llllp{5.2cm}}
    \toprule
    Incident & Khoảng thời gian & Severity & Sự kiện/cảnh báo & Evidence đã ghi nhận \\
    \midrule
    \texttt{incident-36b0fb3f2301} & 11:00--11:00 & Medium & 1 / 1 &
    Tài nguyên hiếm: \texttt{auditor01} truy cập \texttt{accounts}. \\
    \texttt{incident-08588ec64267} & 11:10--11:12 & Critical & 3 / 4 &
    Ba truy cập \texttt{accounts:update} bị từ chối trong cửa sổ 300 giây;
    kèm tín hiệu CTX-REPEATED-DENIAL. \\
    \texttt{incident-38dfcb846bfe} & 22:00--22:05 & Critical & 6 / 10 &
    Ngoài giờ, IP mới, chuỗi năm lần đoán mật khẩu, truy cập \texttt{users:delete}
    bị từ chối và tài nguyên hiếm. \\
    \bottomrule
  \end{tabular}%
  }
\end{table*}


\paragraph{Kiểm tra tính nhất quán.}
Bốn bảng trong bài gồm mô tả dataset, metric baseline, finding ngữ cảnh và
incident; mỗi bảng đều ghi rõ tệp nguồn hoặc artifact CP3 trong caption. Cụ thể,
\texttt{run\_metadata.json} của lần chạy context ghi 14 dòng hợp lệ, 0 dòng lỗi
và 15 alert; các con số 10 finding và 3 incident được đếm từ hai JSON tương
ứng. Vì vậy, so sánh định lượng duy nhất là metric baseline có nhãn; các bảng
ngữ cảnh chỉ mô tả đầu ra của kịch bản tổng hợp và không được dùng để suy ra
precision, recall hay khả năng tổng quát hóa.

\subsection{Case study: chuỗi truy cập rủi ro ngoài giờ}
\label{subsec:case-study}

Incident \texttt{incident-38dfcb846bfe} là ví dụ có nhiều tín hiệu nhất trong
artifact. Bản ghi xác định người dùng \texttt{teller01} và IP
\texttt{198.51.100.50}, với hai mốc 22:00--22:05 ngày 22-06-2026, sáu event và
10 alert. Evidence nguyên gốc gồm truy cập ngoài giờ, IP mới đối với người dùng,
năm lần đăng nhập thất bại trong năm phút, một thao tác
\texttt{users:delete} bị RBAC từ chối và tài nguyên hiếm. Incident nhận điểm
99 và mức Critical, vì chương trình lấy điểm alert lớn nhất trong đoạn và cộng
bonus chuỗi 15 trước khi giới hạn trên 100. Case này minh họa cách các rule nền
và tín hiệu ngữ cảnh tạo một đơn vị triage có thể giải thích, nhưng toàn bộ event
và nhãn là dữ liệu tổng hợp; nó không xác nhận hiệu quả phát hiện trên môi trường
vận hành hoặc hành vi người dùng thật.
